\section{Generated Runtime Classes}
\label{sec:GeneratedRuntime}

Each generated project contains Java classes under
Operational classes (workers, gateways, orchestrator) are generated under:
\filepath{src/main/java/org/gautelis/durga/generated/}.
Probes and exercisers (starter, scenario runner, publishers, watchers) are
generated under:
\filepath{src/main/java/org/gautelis/durga/generated/probes/}.
These classes map BPMN constructs onto Kafka consumers, producers and lifecycle events
rather than hiding behind an opaque engine.

\subsection{Class families}

\begin{itemize}
\item \textbf{Starter classes} publish the initial lifecycle event and first task input.
\item \textbf{Task services} handle BPMN tasks. Service-task handlers auto-complete after
their work step; user and manual task handlers enter a waiting state and await explicit
completion, failure or escalation input.
\item \textbf{Plugin executors} instantiate a registered plugin at field level and delegate
each message to \texttt{plugin.execute()}. No per-plugin code paths in the generator.
\item \textbf{Custom contract interfaces} define the API for hand-written implementations
of custom activities. They extend the \texttt{Plugin} interface.
\item \textbf{Custom delegating workers} inject a contract interface via CDI and delegate
each message to the implementation. Includes a null guard with dead-letter routing.
\item \textbf{Gateway services} handle XOR, AND and OR routing decisions, including fan-out
and join behavior. XOR and OR conditions are evaluated at runtime by a recursive descent
expression parser supporting \texttt{==}, \texttt{!=}, \texttt{>}, \texttt{<}, \texttt{>=},
\texttt{<=}, \texttt{\&\&}, \texttt{||}, \texttt{!}, parentheses, and nested field access.
\item \textbf{Event handlers} implement timer, message and signal behavior for intermediate
events and boundary-event reactions.
\item \textbf{Call-activity helpers} model request and reply behavior for call activities.
\item \textbf{Subprocess scope services} emit subprocess entry and completion lifecycle
events and coordinate subprocess-local routing.
\item \textbf{Process orchestrator} observes terminal flows and emits process completion.
\item \textbf{Probe helpers} provide generated publishers and watchers for command-line
exploration.
\end{itemize}

\subsection{Task services}

Generated task services consume a task input topic, emit \texttt{ACTIVITY\_ENTERED}, perform
the generated action or enter a wait state, then emit \texttt{ACTIVITY\_COMPLETED},
\texttt{FAILED}, \texttt{ESCALATED} or \texttt{CANCELLED} before forwarding to the next
topic or boundary path.

\subsection{Plugin executors and custom workers}

Plugin executors and custom delegating workers share the same structure: a CDI bean
with \texttt{@Incoming} for the task input topic and \texttt{@Channel} emitters for
output, process-events, and dead-letter topics. The difference is in how the
transformation logic is obtained:

\begin{itemize}
\item \textbf{Plugin executor}: instantiates the concrete plugin class as a field
  (\texttt{private final Plugin plugin = new JsonTransform();}) and calls
  \texttt{plugin.execute(payload, config)}.
\item \textbf{Custom delegating worker}: injects the contract interface via CDI
  (\texttt{@Inject MyContract implementation;}) and calls
  \texttt{implementation.execute(payload, config)}. The implementation is provided
  by a developer-written class annotated with \texttt{@ApplicationScoped}.
\end{itemize}

Both patterns include scope cancellation checks, DLQ routing on failure, and
canonical \texttt{process-events} emission.
For both patterns, returned bytes are interpreted as UTF-8. If the result parses
as a JSON object, that object becomes the next \texttt{ProcessEvent.payload}. If
the result is a JSON scalar, a JSON array, or non-JSON text, the generated worker
wraps it as \texttt{\{ "\_value": ... \}}. A \texttt{null} result keeps the incoming
payload unchanged.

\subsection{Routing and event classes}

Gateways turn BPMN branching into explicit topic routing. Timer, message and signal handlers
make delays and external triggers concrete. Boundary handlers observe task or scope state and
branch into exceptional or reminder flows when trigger conditions are met.

\subsection{Scope classes}

Subprocess scope services emit subprocess-level entered and completed events, mediate flow
into and out of the subprocess graph, propagate failure, escalation and cancellation at
scope level, and support nested subprocess structure. Event subprocess handlers model
subprocess-internal event-driven branches separately from the main task path.
